%% Where the parentheses live: what \ref prints under [langsci], and what
%% does not follow it.
%%
%% langsci-gb4e leaves the counter bare and parenthesises the LIST LABEL,
%% so \ref there gives "1" and an author writes "(\ref{ex:x})" or \xref;
%% linguex puts the parentheses in \theExNo, so \ref gives "(1)".  A
%% document ported from the first to the second prints "((1))" until the
%% flavour follows the syntax, which is what [langsci] now makes it do.
%%
%% The other half of the rule is the half a one-sided fix breaks: every
%% number this package prints ITSELF is parenthesised whatever the flavour
%% -- the example's own label, \Next and \Last, \xref, \xxref, \exr and
%% \Refrange -- so the same page must show bare and parenthesised numbers
%% side by side.  A change that took the parentheses off \theExNo alone
%% would strip the label with them and still pass a case that only looked
%% at \ref.
%%
%% \ExParenRefs and \ExBareRefs ask for the other flavour by name, and are
%% here in both directions: the second half of this case is a [langsci]
%% document put back to linguex references and then back again.
\input{_preamble-langsci}
\begin{document}
\ea\label{lr:main} LRMAIN a main example.
\z

\ea LRHOST a host.
  \ea\label{lr:sub} LRSUBA a sub-example.
  \ex LRSUBB another.
    \ea\label{lr:rom} LRROMA a roman one.
    \z
  \z
\z

LRREF \ref{lr:main} \ref{lr:sub} \ref{lr:rom} LRREFEND

LRXREF \xref{lr:main} LRXXREF \xxref{lr:main}{lr:sub} LRRANGE
\Refrange{lr:main}{lr:sub} LRREL \Next\ \Last

%% The footnote series references the same way, on its own numbering.
LRFNHOST a host sentence.\footnote{\ea\label{lr:fn} LRFN an example in a
  footnote.\z LRREFFN \ref{lr:fn} LRREFFNEND}

\ExParenRefs
\ea\label{lr:paren} LRPAREN under \string\ExParenRefs.
\z
LRREFPAREN \ref{lr:paren} LRREFPARENEND

\ExBareRefs
\ea\label{lr:bare} LRBARE back under \string\ExBareRefs.
\z
LRREFBARE \ref{lr:bare} LRREFBAREEND
\end{document}
